% Ch. 36 : MapReduce, compter les mots des titres (map, combineur, shuffle, reduce)
\documentclass[border=6pt]{standalone}
\usepackage{coursfig}
\begin{document}
\begin{tikzpicture}[x=1mm,y=1mm]
  % colonnes : morceaux, map, combineur, reduce, résultat
  \node[ftitle] at (0,14) {morceaux};
  \node[ftitle] at (46,14) {map};
  \node[ftitle] at (92,14) {combineur};
  \node[ftitle] at (148,14) {reduce};
  \node[ftitle] at (196,14) {résultat};
  \foreach \i/\y/\t/\p in {1/0/{payer le loyer\\appeler le médecin}/{(payer,1) (le,1)\\(loyer,1) (appeler,1)\\(le,1) (médecin,1)},
                           2/-26/{ranger la cave\\laver la voiture}/{(ranger,1) (la,1)\\(cave,1) (laver,1)\\(la,1) (voiture,1)},
                           3/-52/{acheter le pain\\sortir la poubelle}/{(acheter,1) (le,1)\\(pain,1) (sortir,1)\\(la,1) (poubelle,1)}} {
    \node[box=Rule, fill=white, text width=30mm, minimum height=16mm, font=\sffamily\normalsize, align=center] (m\i) at (0,\y) {\t};
    \node[box=Sea, text width=34mm, minimum height=16mm, font=\ttfamily\small, align=center] (p\i) at (46,\y) {\p};
    \draw[flow=Sea] (m\i) -- (p\i);
  }
  \node[box=Ocre, text width=30mm, minimum height=16mm, font=\ttfamily\small, align=center] (c1) at (92,0) {(le,2) (payer,1)\\(loyer,1) \dots};
  \node[box=Ocre, text width=30mm, minimum height=16mm, font=\ttfamily\small, align=center] (c2) at (92,-26) {(la,2) (ranger,1)\\(cave,1) \dots};
  \node[box=Ocre, text width=30mm, minimum height=16mm, font=\ttfamily\small, align=center] (c3) at (92,-52) {(le,1) (la,1)\\(pain,1) \dots};
  \foreach \i in {1,2,3} { \draw[flow=Ocre] (p\i) -- (c\i); }
  % reducers
  \node[box=Enc, text width=30mm, minimum height=18mm, align=center] (r1) at (148,-8) {\textbf{reducer 1}\sub{mots dont\\crc32 mod 2 = 0}};
  \node[box=Enc, text width=30mm, minimum height=18mm, align=center] (r2) at (148,-44) {\textbf{reducer 2}\sub{mots dont\\crc32 mod 2 = 1}};
  \foreach \i in {1,2,3} { \foreach \r in {1,2} { \draw[flow=Plum] (c\i.east) -- (r\r.west); } }
  \node[elabel, text=Plum, font=\sffamily\normalsize\bfseries] at (122,-26) {shuffle};
  \node[box=Rule, fill=white, text width=26mm, minimum height=18mm, font=\ttfamily\small, align=center] (o1) at (196,-8) {(le,3)\\(payer,1) \dots};
  \node[box=Rule, fill=white, text width=26mm, minimum height=18mm, font=\ttfamily\small, align=center] (o2) at (196,-44) {(la,3)\\(cave,1) \dots};
  \draw[flow=Enc] (r1) -- (o1);
  \draw[flow=Enc] (r2) -- (o2);
  \draw[brace] (26,-66) -- node[below=2mm, font=\sffamily\normalsize, text=Muted]{sur chaque machine, sans rien échanger} (-16,-66);
  \draw[brace] (212,-66) -- node[below=2mm, font=\sffamily\normalsize, text=Muted, align=center]{chaque mot arrive chez un seul reducer} (132,-66);
  \node[note, anchor=north west, text width=226mm] at (-16,-78) {mesuré sur 5 millions de titres (8 mappers, 4 reducers) : sans combineur, 22\,680\,376 paires traversent le shuffle, et le reducer le plus chargé en reçoit 8,5 millions contre 3,6 millions au moins chargé (« le » et « la » tombent chez lui) ; avec combineur, il n'en passe plus que 1\,968};
\end{tikzpicture}
\end{document}
